\chapter{Models, mida i comparació amb \texorpdfstring{ZF\slash ETCS}{ZF/ETCS}}\label{cap:models}

\begin{objectius}
\apartat{Prerequisits} Tot el llibre, i en particular els axiomes tal com s'han anat introduint, l'esquema SR del capítol~\ref{cap:alephs} i les observacions de lectura clàssica. Familiaritat informal amb els axiomes de ZFC, que es recorden a la taula de la secció~\ref{sec:taula-comparativa}.
\apartat{Objectius} En acabar el capítol, el lector hauria de poder:
\begin{itemize}
  \item recórrer els axiomes del llibre per grups i explicar-ne les dependències;
  \item identificar l'estructura obtinguda com un topos ben apuntat amb objecte dels nombres naturals i elecció, és a dir, un model d'ETCS;
  \item comparar, axioma per axioma, aquesta presentació amb ZFC i amb l'ETCS de Lawvere, distingint els axiomes equivalents, els derivats i els addicionals;
  \item situar la força de la teoria (Nivell~I, amb i sense SR) respecte de la teoria de Zermelo acotada i de ZFC;
  \item explicar com es recupera la pertinença, què se'n perd i què se'n guanya.
\end{itemize}
\end{objectius}

Els capítols anteriors han bastit, a partir d'un llenguatge en què els
únics individus són les aplicacions, tota la infraestructura que
habitualment associem als conjunts: parts i pertinença, productes i sumes,
quocients, imatges, exponencials, el conjunt de parts, els nombres
naturals i la cardinalitat fins al numerable. I, en una segona part, han
portat aquesta construcció fins a la teoria clàssica de l'infinit: la
comparació general de mides amb Schröder--Bernstein i el teorema de Cantor
(capítol~\ref{cap:cantor}), els bons ordres i la recursió transfinita
(capítols~\ref{cap:bons-ordres} i~\ref{cap:recursio}), Zorn, Hartogs i el
teorema del bon ordre (capítol~\ref{cap:zorn}), els tipus ordinals i la seva
aritmètica (capítols~\ref{cap:ordinals} i~\ref{cap:aritmetica-ordinal}), els
cardinals i la seva aritmètica (capítol~\ref{cap:cardinals}) i ---amb un principi
addicional--- la jerarquia sencera dels alephs i dels beths
(capítol~\ref{cap:alephs}). En cap
moment hem necessitat la relació de pertinença \(\in\) com a noció primitiva.
És natural que, un cop acabada la construcció, el lector es faci tres
preguntes: si això que hem obtingut és \emph{realment} una teoria de
conjunts; què és capaç d'expressar i què no; i quina relació guarda amb la
teoria clàssica de Zermelo--Fraenkel que va aprendre primer. Aquest capítol
final hi respon. No conté cap axioma nou ni cap construcció nova: és un
capítol metateòric, escrit per situar l'edifici sencer ---inclosa la part
transfinita--- i per assenyalar-ne els límits amb honestedat.

\section{Recapitulació dels axiomes}\label{sec:recapitulacio-axiomes}

Convé, abans de res, tenir a la vista tots els axiomes que hem assumit i
l'ordre en què depenen els uns dels altres. Els agrupem pel paper que fan.

El \textbf{grup del llenguatge de composició} fixa què és una aplicació i
com es componen: l'existència i unicitat de domini i codomini
(axioma~\ref{ax:dom-cod}), l'existència de la composició quan els tipus
encaixen (axioma~\ref{ax:composicio-existencia}), el seu domini i codomini
(axioma~\ref{ax:composicio-dom-cod}), l'associativitat
(axioma~\ref{ax:associativitat}) i les lleis de neutralitat
(axioma~\ref{ax:neutralitat}). Aquests cinc axiomes són la base sobre la
qual descansa tota la resta; en el vocabulari que hem evitat, diuen que les
aplicacions i els conjunts formen una categoria.

El \textbf{grup dels objectes distingits i la separació} dona els punts de
partida per construir: l'existència d'un conjunt terminal
(axioma~\ref{ax:terminal-existencia}) i d'un d'inicial
(axioma~\ref{ax:inicial-existencia}), la no-degeneració que els separa,
\(\mathbf0\ncong\mathbf1\) (axioma~\ref{ax:no-degeneracio}), i la separació
de les aplicacions pels elements globals
(axioma~\ref{ax:separacio-globals}), que és la que fa que els punts
\(\mathbf1\to A\) siguin prou nombrosos per distingir aplicacions ---la
propietat que més endavant expressarem dient que la teoria és \emph{ben
apuntada}.

El \textbf{grup de les construccions finites} postula les operacions
universals bàsiques: l'existència de sumes (axioma~\ref{ax:sumes}), la seva
disjunció i estabilitat (axioma~\ref{ax:sumes-disjuntes}), l'existència de
productes (axioma~\ref{ax:productes}), la representació de les parts per
aplicacions cap a \(\mathbf2\) (axioma~\ref{ax:representacio}) ---el nucli
lògic de tot el llibre--- i l'existència de coequalitzadors
(axioma~\ref{ax:coequalitzadors}). Amb ells s'obtenen tots els límits i
colímits finits i tota l'àlgebra booleana de parts.

El \textbf{grup de la regularitat i l'exactitud} controla les imatges i
els quocients: l'efectivitat de les relacions d'equivalència
(axioma~\ref{ax:exactitud}), que casa quocients i parells nucli. Combinat
amb els coequalitzadors i la separació global, dona la factorització
imatge, i amb ella que tot epimorfisme és regular i estable per canvi de
base, sense elecció (capítol~\ref{cap:imatges}).

El \textbf{grup de l'ordre superior} fa la teoria capaç de parlar de si
mateixa: l'existència d'exponencials (axioma~\ref{ax:exponencials}), de la
qual surt, prenent \(\mathbf2^{A}\), el conjunt de parts com a objecte, la
pertinença interna i els quantificadors.

El \textbf{grup de l'infinit} afegeix l'única font d'infinitud: l'objecte
dels nombres naturals amb la seva propietat universal de recursió
(axioma~\ref{ax:naturals}), del qual pengen la inducció, l'aritmètica, la
finitud interna i tota la cardinalitat numerable.

Finalment, els \textbf{axiomes d'elecció}, que hem mantingut separats
perquè molts resultats no en depenen: l'elecció global ---tot epimorfisme
s'escindeix--- (axioma~\ref{ax:eleccio}) i, com a cas particular seu,
l'elecció numerable (axioma~\ref{ax:eleccio-numerable}), que va ser
necessària per identificar les dues nocions d'infinitud
(teorema~\ref{teo:infinit-implica-dedekind}) i per la unió numerable
d'enumerables (teorema~\ref{teo:unio-numerable}).

Fora d'aquesta llista, i amb un estatus deliberadament diferent, hi ha
l'\textbf{esquema de reunió de famílies} o \emph{substitució estructural}
(principi~\ref{sr:esquema}), que la part transfinita ha necessitat només al
capítol opcional dels alephs (capítol~\ref{cap:alephs}). Insistim que SR \emph{no} és un axioma més al costat dels anteriors: és un
\emph{esquema} ---una instància per a cada fórmula funcional i invariant per
isomorfisme---, no una sentència, exactament com el reemplaçament de la
teoria clàssica és un esquema i no un axioma únic. Per això l'hem presentat
com un \emph{principi} a part i no amb l'entorn dels axiomes, i per això tots
els capítols fins al~18 en són independents: la teoria de Nivell~I ---tota
la que va del llenguatge d'aplicacions fins als cardinals--- no el fa servir.
Només la reunió de famílies transfinites no acotades ---la successió dels
alephs, els beths--- el requereix. Aquesta és la línia divisòria que
recorrerem amb detall a la secció~\ref{sec:etcs-sr}.

\begin{observacio}[El mapa de dependències]\label{obs:mapa-dependencies}
La lectura del llibre segueix, de fet, aquest ordre de dependència. El
llenguatge de composició no pressuposa res; els objectes distingits el
necessiten; les construccions finites necessiten els objectes distingits;
la regularitat necessita sumes i coequalitzadors; els exponencials
necessiten els productes; l'objecte dels nombres naturals necessita els
exponencials per a la recursió amb paràmetres
(observació~\ref{obs:exponencials-recursio}); i l'elecció és transversal,
invocada només on s'assenyala. Cap axioma no es fa servir abans
d'enunciar-se, i cap grup no depèn d'un de posterior. Aquesta és la raó
per la qual el llibre es pot llegir linealment: l'estructura lògica i
l'estructura expositiva coincideixen.
\end{observacio}

\section{Què és, categòricament, l'estructura obtinguda}\label{sec:caracteritzacio}

Reunits tots els axiomes, es pot dir en una sola frase quina mena
d'estructura formen les aplicacions i els conjunts d'aquest llibre. En el
vocabulari de la teoria de categories ---que no hem fet servir en cap
demostració, però que aquí, en un capítol de balanç, resulta el més
econòmic per anomenar el conjunt---, l'estructura és un \emph{topos ben
apuntat, amb objecte dels nombres naturals i elecció}. Desglossem-ho amb
paraules planes, perquè cada peça correspon exactament a un grup
d'axiomes.

Ser un \emph{topos elemental} vol dir tres coses, que són les que hem
postulat: que hi ha tots els límits finits ---terminal, productes,
equalitzadors i les seves combinacions, que vam obtenir al
capítol~\ref{cap:valors}---; que hi ha exponencials
(axioma~\ref{ax:exponencials}); i que les parts d'un conjunt qualsevol
estan representades per les aplicacions cap a un objecte fix \(\mathbf2\),
el \emph{classificador de subobjectes} (axioma~\ref{ax:representacio}).
Aquesta darrera és la propietat decisiva: fa que la lògica dels
subconjunts sigui interna a la teoria, amb el conjunt de parts, la
pertinença i els quantificadors com a objectes i aplicacions de ple dret,
tal com vam veure al capítol~\ref{cap:parts-interna}.

\emph{Ben apuntat} (\emph{well-pointed}) és el que afegeix l'axioma de separació pels elements
globals (axioma~\ref{ax:separacio-globals}) juntament amb la
no-degeneració (axioma~\ref{ax:no-degeneracio}): els punts
\(\mathbf1\to A\) basten per distingir aplicacions, i \(\mathbf1\) no
coincideix amb \(\mathbf0\). És el que permet raonar «element a element»,
com fem contínuament amb els elements globals i generalitzats, sense
trair el punt de vista estructural. Un topos qualsevol no té per què ser
ben apuntat; el nostre sí, i això l'acosta molt a la categoria dels
conjunts ordinaris.

L'\emph{objecte dels nombres naturals} (axioma~\ref{ax:naturals}) i
l'\emph{elecció} (axioma~\ref{ax:eleccio}) són els dos ingredients que
completen el quadre. Un topos ben apuntat amb aquests dos afegits és,
precisament, un model de la teoria estructural de conjunts de Lawvere,
l'ETCS~\cite{lawvere1964,lawvere-etcs-reprint}. Dit d'una altra manera:

\begin{observacio}[La tesi del llibre, en una frase]\label{obs:tesi-etcs}
Els axiomes d'aquest llibre són una axiomatització d'ETCS presentada en un
llenguatge de primer ordre d'aplicacions, sense el vocabulari de la teoria
de categories. Tot el que hem demostrat és, doncs, un teorema d'ETCS; i
qualsevol model d'ETCS ---en particular la categoria dels conjunts
clàssics--- satisfà els nostres axiomes. Aquesta és la resposta curta a la
pregunta «és realment una teoria de conjunts?»: ho és tant com ETCS, que
és una axiomatització \emph{autònoma} de la categoria dels conjunts,
proposada com a fonament alternatiu de la matemàtica~\cite{lawvere-rosebrugh,mclarty-catset}.
Evitem l'expressió \enquote{axiomatització completa}, que es podria confondre
amb la \emph{completesa lògica} (que tot enunciat o la seva negació és
demostrable) o amb la \emph{categoricitat} (que hi ha un únic model llevat
d'isomorfisme): cap de les dues no val aquí ---ETCS és una teoria de primer
ordre, subjecta a la incompletesa i amb molts models no isomorfs---, ni cap de
les dues no es demostra en aquest llibre. \enquote{Autònoma} vol dir només que
els axiomes es basten a si mateixos com a presentació dels conjunts, sense dependre d'una teoria material prèvia.
\end{observacio}

\section{La part transfinita dins el balanç}\label{sec:transfinit-balanc}

El capítol que en la primera versió d'aquest llibre tancava el volum
s'aturava a la cardinalitat numerable. La part transfinita
(capítols~\ref{cap:cantor}--\ref{cap:alephs}) ha estès el desenvolupament
fins a la teoria clàssica de l'infinit, i convé situar-la dins aquest
balanç, perquè il·lumina exactament què es pot fer amb els axiomes de
Nivell~I i què demana un principi addicional.

El primer bloc és \emph{purament diagonal} i no costa cap axioma nou. El
teorema de \textbf{Schröder--Bernstein} (teorema~\ref{teo:schroder-bernstein})
---dues injeccions oposades donen un isomorfisme--- es prova pel menor punt
fix d'un operador creixent sobre \(P(A)\), sense elecció. El teorema de
\textbf{Cantor} (teorema~\ref{teo:cantor}) ---\(A\prec P(A)\)--- és
l'argument diagonal, i tampoc no necessita res de nou; d'ell surt que no hi
ha conjunt universal ni mida màxima. Aquests dos resultats són, en el
llenguatge estructural, teoremes elementals sobre parts i pertinença interna.

El segon bloc és la \emph{teoria de l'ordre}, també de Nivell~I. Els
\textbf{bons ordres} (capítol~\ref{cap:bons-ordres}) tenen la rigidesa que
els fa mesura de forma ---entre dos n'hi ha com a molt un isomorfisme
d'ordre---, i sobre ells es construeixen la \textbf{inducció} i la
\textbf{recursió transfinites} amb codomini fix (capítol~\ref{cap:recursio}),
la reunió de solucions parcials dins un exponencial fix, \emph{sense}
elecció. Amb la recursió es tanca la comparabilitat de bons ordres
(teorema~\ref{teo:comparabilitat}). El teorema de \textbf{Hartogs}
(teorema~\ref{teo:hartogs}) ---un bon ordre que no s'injecta en \(A\)--- és
el paral·lel ordinal de Cantor, i és igualment de Nivell~I: es construeix
tallant els bons ordres sobre parts d'\(A\) ---parells \((S,R)\) dins
\(P(A)\times P(A\times A)\)--- i quocientant-los per isomorfisme.

El tercer bloc introdueix l'\emph{elecció}. El teorema del \textbf{bon
ordre} (teorema~\ref{teo:bon-ordre}) i el \textbf{lema de Zorn}
(capítol~\ref{cap:zorn}) es demostren amb l'elecció i la cota de Hartogs
---que evita, precisament, la recursió \enquote{sobre tots els ordinals}---,
i resulten equivalents a l'elecció. Sobre aquesta base, els \textbf{tipus
ordinals} i la seva aritmètica (capítols~\ref{cap:ordinals}
i~\ref{cap:aritmetica-ordinal}) i els \textbf{cardinals} i la seva
(capítol~\ref{cap:cardinals}) es desenvolupen amb el teorema de Cantor
reaparegut com \(\abs{A}<2^{\abs{A}}\), la comparabilitat total dels
cardinals i el teorema de Burali--Forti (teorema~\ref{teo:burali-forti}) com
a segon gran resultat diagonal ---no hi ha longitud d'ordre màxima.

Tot això és de Nivell~I o de Nivell~I més elecció: viu dins conjunts ja
construïts, acotat per un \(\Hart\) sempre que cal reunir tipus. \emph{Només}
el capítol final dels alephs (capítol~\ref{cap:alephs}) surt d'aquest terreny,
i ho fa exactament en un punt: reunir la successió transfinita no acotada
\(\alpha\mapsto\aleph_\alpha\) ---i la dels beths--- en una família. Aquest
és l'únic lloc de tot el llibre on cal SR, i és el tema de la
secció~\ref{sec:etcs-sr}.

\section{Ordinal material i tipus estructural}\label{sec:vonneumann}

La part transfinita permet, a més, precisar una comparació que la
introducció només podia insinuar: la que hi ha entre l'\emph{ordinal
material} de la teoria clàssica i el \emph{tipus d'ordre estructural}
d'aquest llibre.

En la teoria de von Neumann, un ordinal \emph{és} un conjunt: un conjunt
transitiu ben ordenat per la pertinença, de manera que cada ordinal és
literalment el conjunt dels ordinals menors ---\(0=\varnothing\),
\(1=\{0\}\), \(2=\{0,1\}\), \(\omega=\{0,1,2,\dots\}\)--- i
\enquote{\(\alpha<\beta\)} es defineix com \enquote{\(\alpha\in\beta\)}. És
una construcció elegant i canònica: cada tipus de bon ordre té \emph{un}
representant distingit, l'ordinal de von Neumann, i la relació d'ordre és la
pertinença mateixa. El preu és que descansa sobre la pertinença global i la
jerarquia acumulativa: un ordinal arrossega tota la seva història de
membres, i la seva construcció ---per recursió al llarg de tots els
ordinals--- fa servir el reemplaçament.

En aquest llibre un \emph{tipus ordinal} no és cap conjunt privilegiat, sinó
un bon ordre \emph{considerat llevat d'isomorfisme d'ordre}
(definició~\ref{def:tipus-ordinal}), i \enquote{\(\alpha<\beta\)} es defineix
per segments inicials, no per pertinença ---de fet vam advertir explícitament
que \(\alpha\in\beta\) no és una notació adequada. Els dos enfocaments
descriuen la mateixa matemàtica, i la traducció és neta: el tipus estructural
d'un bon ordre \(W\) correspon a la classe d'isomorfisme d'ordre de \(W\),
i l'ordinal material de von Neumann n'és el representant canònic que la teoria
clàssica \emph{tria} dins aquesta classe. On l'estructural diu \enquote{el
tipus de \(W\), únic llevat d'un únic isomorfisme}, el material diu
\enquote{l'ordinal \(\in\)-transitiu isomorf a \(W\)}. La diferència no és de
contingut sinó de si es fa o no una tria canònica de representant: el
tractament estructural la evita perquè no la necessita, i perquè triar-la
demanaria la mateixa pertinença global que hem volgut no fer primitiva.

\begin{observacio}[Per què l'estructural no tria representant]\label{obs:vonneumann-representant}
Es podria pensar que l'ordinal de von Neumann és \enquote{millor} perquè és
canònic. Però la seva canonicitat és exactament el que costa la pertinença
global: per triar, dins cada classe d'isomorfisme de bons ordres, l'ordinal
\(\in\)-transitiu, cal disposar de la relació \(\in\) i de la recursió que el
construeix ---i, per als ordinals grans, del reemplaçament. El tractament
estructural prefereix treballar amb el tipus \emph{sense} representant, com fa
la pràctica matemàtica quan diu \enquote{el tipus d'ordre de \(\Q\)} sense
identificar-lo amb cap conjunt concret. Quan es vol un representant ---per
exemple per fer aritmètica còmoda--- se'n pot triar un amb l'elecció (el
cardinal inicial n'és un exemple, capítol~\ref{cap:cardinals}), però és una
tria \emph{addicional}, no un tret primitiu. Aquesta és la mateixa filosofia
que separa \(\abs{A}\) com a classe d'equipotència del cardinal inicial
escollit: l'estructura reté el que és invariant, i deixa la tria de
representant per a qui la vulgui.
\end{observacio}

\section{Interpretació en conjunts petits}\label{sec:conjunts-petits}

Una manera de convèncer-se que els axiomes descriuen conjunts és
interpretar-los en un univers ja conegut i comprovar que s'hi compleixen.
Fixem, doncs, un univers de conjunts clàssics ---els conjunts d'una teoria
de Zermelo--Fraenkel, posem--- i llegim-hi cada noció del llibre.

Un \emph{conjunt} del llibre és, en aquesta lectura, un conjunt ordinari.
Una \emph{aplicació} és una funció en el sentit habitual, amb el seu
domini i el seu codomini fixats ---no un simple graf, sinó una funció amb
els seus dos extrems marcats, que és el que demana l'axioma~\ref{ax:dom-cod}.
La \emph{composició} és la composició de funcions; les \emph{identitats}
són les funcions identitat. El conjunt \emph{terminal} \(\mathbf1\) és
qualsevol singletó \(\{\ast\}\); l'\emph{inicial} \(\mathbf0\) és el
conjunt buit; la \emph{no-degeneració} és el fet que \(\varnothing\) i
\(\{\ast\}\) no són equipotents. Un \emph{element global} \(\mathbf1\to A\)
és, exactament, un element de \(A\) en el sentit ordinari, perquè triar
una funció des d'un singletó és triar un valor. El conjunt \(\mathbf2\) és
qualsevol parell \(\{0,1\}\), i una \emph{part} de \(A\) es correspon amb
un subconjunt ordinari via la seva funció característica
(axioma~\ref{ax:representacio}). Les \emph{sumes} són unions disjuntes; els
\emph{productes}, productes cartesians; els \emph{coequalitzadors},
quocients per relacions d'equivalència; els \emph{exponencials}
\(B^{A}\), el conjunt de totes les funcions de \(A\) a \(B\); l'\emph{objecte
dels nombres naturals}, el conjunt \(\{0,1,2,\dots\}\) amb el zero i el
successor.

\begin{observacio}[La lectura és fidel, però no és l'única]\label{obs:interpretacio-fidel}
Sota aquesta interpretació, cadascun dels divuit axiomes ---els disset
estructurals més l'elecció global, sense comptar-hi l'elecció numerable,
que n'és un cas particular--- és un teorema
elemental de la teoria clàssica de conjunts: que tota funció té domini i
codomini, que el producte cartesià existeix, que el conjunt de funcions
existeix, i així successivament. Per això la categoria dels conjunts
clàssics és un model dels nostres axiomes, i qualsevol cosa que hàgim
demostrat hi val. Ara bé ---i això és essencial---, \emph{no és l'únic}
model. Els mateixos axiomes se satisfan en universos molt més petits o de
mides diferents: per exemple, la col·lecció dels conjunts \emph{finits}, amb
les funcions entre ells, compleix tots els axiomes \emph{llevat} del de
l'infinit (axioma~\ref{ax:naturals}), i restringint-se a un univers de
Grothendieck (un univers de conjunts de rang menor que un cardinal
inaccessible) s'obtenen models complets tancats sota totes les operacions,
inclòs el conjunt de parts. La teoria no fixa un univers; en descriu una
\emph{forma}, i molts universos de mides diferents la tenen.
\end{observacio}

\begin{observacio}[Per què no serveixen els hereditàriament numerables]\label{obs:no-hcountable}
Podria temptar prendre com a model \enquote{petit} la col·lecció dels conjunts
\emph{hereditàriament numerables}. No funciona: aquesta col·lecció \emph{no}
és tancada sota el conjunt de parts. En efecte, \(\N\) és hereditàriament
numerable, però \(P(\N)\) no ho és ---és no numerable---, de manera que
l'axioma del conjunt de parts falla dins d'aquell univers. Un model complet
dels axiomes ha de ser tancat sota totes les operacions que aquests postulen;
per això cal, o bé un univers de Grothendieck (tancat sota parts amb el marge
de mida corresponent), o bé un univers amb hipòtesis de mida explícites que
garanteixin la clausura. La numerabilitat hereditària, tota sola, no la dona.
\end{observacio}

Aquesta pluralitat de models no és un defecte sinó el senyal distintiu
d'una axiomatització estructural: descriu els conjunts \emph{llevat
d'isomorfisme i llevat d'equivalència de la categoria sencera}, sense
comprometre's amb cap realització particular. És el mateix que passa amb
els axiomes de grup, que no descriuen \emph{un} grup sinó la forma comuna
a tots.

\section{Pertinença primitiva i pertinença recuperada}\label{sec:pertinenca}

La diferència més visible amb la teoria clàssica és el paper de la
pertinença. Val la pena precisar-la, perquè és on es concentra bona part
del canvi de perspectiva anunciat a la introducció.

En la teoria de Zermelo--Fraenkel, \(\in\) és una relació \emph{primitiva}
i \emph{global}: té sentit preguntar-se, per a dos conjunts qualssevol
\(x\) i \(A\), si \(x\in A\), i la resposta forma part de les dades del
món. Que \(3\in 5\) (en la codificació de von Neumann els naturals són
conjunts, i \(5=\{0,1,2,3,4\}\)) és un fet; que \(3\in\pi\), amb una certa
codificació dels reals, també ho és, encara que sigui un fet estrany.
Aquesta globalitat té una conseqüència: cada conjunt arrossega tota la
seva història interna ---els seus elements, els elements dels seus
elements, i així avall fins al buit---, i l'axioma de fonamentació
garanteix que aquesta baixada acaba.

En la teoria d'aquest llibre no hi ha cap relació \(\in\) primitiva. El
que hi ha és una pertinença \emph{recuperada}, i \emph{local}. «Ser
element de \(A\)» vol dir ser un element global \(\mathbf1\to A\), o més
generalment un element generalitzat \(X\to A\); i «pertànyer a una part
\(P\hookrightarrow A\)» vol dir factoritzar per la inclusió, cosa que la
pertinença interna del capítol~\ref{cap:parts-interna} converteix en una
aplicació \(A\to\mathbf2\) avaluada al punt. Aquesta pertinença sempre és
\emph{relativa a un conjunt ambient}: un element ho és \emph{d'un} conjunt
determinat, i la pregunta «aquest element pertany a aquell altre conjunt?»
només té sentit un cop fixat el codomini.

\begin{observacio}[Què es perd i què es guanya]\label{obs:pertinenca-local}
La pèrdua és la pertinença creuada entre conjunts arbitraris: aquí no té
sentit preguntar si \(\N\) és element de \(\R\), ni si un element de \(A\)
«és el mateix» que un element de \(B\) quan \(A\) i \(B\) no tenen cap
relació ---els elements no tenen identitat pròpia fora del conjunt que els
allotja. En la lectura estructural això no és una mancança sinó una
disciplina: reflecteix que en la pràctica matemàtica mai no comparem un
nombre natural amb un nombre real «com a conjunts», sinó a través d'una
aplicació explícita \(\N\to\R\). El guany és que desapareixen els fets
accidentals de codificació ---el \(3\in\pi\) d'abans--- que la teoria
clàssica genera i que cap matemàtic no fa servir. La informació que queda
és exactament la invariant per isomorfisme, que és la que la matemàtica
utilitza.
\end{observacio}

\begin{observacio}[Es pot recuperar una \(\in\) global?]\label{obs:recuperar-in}
Sí, i el fet que es pugui és el pont tècnic amb la teoria clàssica. A
partir d'un model dels nostres axiomes ---amb un afegit sobre la mida que
esmentarem--- es pot construir una relació \(\in\) global i un univers de
conjunts «amb elements», recuperant un model de (una part de) Zermelo. La
construcció associa a cada conjunt estructural un arbre d'aplicacions que
en codifica una versió «ben fundada», i és inversa, en el sentit adequat,
de la lectura de la secció~\ref{sec:conjunts-petits}. No la desenvolupem
aquí; el que importa és que la pertinença global no és una dada primitiva
irrenunciable, sinó una estructura \emph{derivable} quan es vol, a canvi
d'hipòtesis addicionals de reemplaçament~\cite{mclarty-catset,shulman-comparing}.
\end{observacio}

\section{Comparació amb ZF/ZFC i amb ETCS}\label{sec:comparacio}

Amb això a la mà, podem situar el llibre respecte de les dues grans
famílies d'axiomatitzacions.

Respecte de \textbf{ZF/ZFC}, la diferència de fons no és de potència sinó
de \emph{primitius}. ZF pren com a primitius els conjunts i la relació
\(\in\), i construeix les funcions com a conjunts de parells; nosaltres
prenem com a primitives les aplicacions, i recuperem la pertinença com una
noció derivada i local, tal com acabem de veure (secció~\ref{sec:pertinenca}).
La diferència no és que ZF descrigui un univers únic i la nostra teoria molts:
ZF és igualment una teoria de primer ordre amb molts models no isomorfs (pel
teorema de Löwenheim--Skolem i la incompletesa, en té de mides i tipus
d'isomorfisme diferents). El contrast correcte és de \emph{criteri d'identitat}:
en la lectura material de ZF, cada element està determinat \emph{absolutament}
per la relació \(\in\) ---dos conjunts amb els mateixos elements són el mateix,
i la pertinença fixa la identitat de cada objecte dins un model donat---,
mentre que en la lectura estructural la identitat és \emph{llevat
d'isomorfisme}: un objecte només està determinat pel seu paper en la categoria,
i cap propietat interna distingeix objectes isomorfs. La distinció, doncs, és
entre identitat rígida per \(\in\) i identitat estructural llevat
d'isomorfisme, no entre una teoria amb model únic i una amb molts models. En el
pla lògic, totes dues són teories de primer ordre, però amb signatures
diferents: la de ZF té un únic símbol de relació \(\in\); la nostra, els
predicats \(\Dom\), \(\Cod\) i \(\Com\) sobre aplicacions.

Respecte d'\textbf{ETCS}, la relació és de \emph{presentació}, no de
contingut. Com ja hem dit (observació~\ref{obs:tesi-etcs}), els nostres
axiomes són ETCS. La diferència amb l'exposició de Lawvere és que ell la
formula en el llenguatge de les categories ---parla de límits, objectes
exponencials i classificador de subobjectes--- mentre que nosaltres hem
volgut que cada axioma fos una fórmula elemental sobre aplicacions,
llegible sense saber teoria de categories, i que les construccions
categòriques apareguessin com a \emph{teoremes} i no com a vocabulari de
partida. El contingut matemàtic és el mateix; canvia qui és el públic.

\begin{observacio}[Força relativa: el lloc exacte a l'escala]\label{obs:forca-relativa}
La comparació de força mereix una mica de cura, perquè el resultat sovint
sorprèn. ETCS ---i per tant aquest llibre--- \emph{no} és tan fort com
ZFC. La seva força demostrativa és la de la teoria de Zermelo amb elecció
i amb separació restringida a fórmules \emph{acotades} (l'anomenada
\(Z\) acotada, o \(\mathrm{BZC}\)): ETCS i \(\mathrm{BZC}\) demostren
exactament els mateixos enunciats aritmètics, són \emph{equiconsistents} i
mútuament interpretables~\cite{osius1974,mclarty-catset}. El que separa ETCS de ZFC
és l'\emph{esquema de reemplaçament}, que ETCS en la seva forma bàsica no
inclou. Afegint a ETCS una versió categòrica concreta del reemplaçament ---la
formulació de reemplaçament\slash col·lecció estructural adoptada com a SR
(observació~\ref{obs:sr-rigor}, capítol~\ref{cap:alephs})--- s'obté una teoria
de la mateixa força que ZFC, i llavors la correspondència amb la teoria
clàssica és exacta en el sentit precís de la bi-interpretabilitat demostrada
a~\cite{osius1974,shulman-comparing} (no en el sentit de completesa lògica).
La correspondència depèn, doncs, de \emph{quina} versió de l'esquema
s'adopti, i és per a aquella versió que val el teorema de comparació. Així, en
l'escala de les teories de conjunts, aquest llibre se situa una mica per
sota de ZFC ---al nivell de la teoria de Zermelo amb elecció---, que és
més que suficient per a pràcticament tota la matemàtica ordinària, però
que cal saber reconèixer per no atribuir a la teoria una potència que no
té sense l'afegit del reemplaçament.
\end{observacio}

\subsection{ETCS, reemplaçament i la substitució estructural}\label{sec:etcs-sr}

La discussió anterior sobre la força relativa pren un relleu concret ara que
el llibre ha desenvolupat la part transfinita, perquè és allà on la
diferència entre ETCS i ETCS reforçat es fa tangible. Convé, doncs,
distingir amb precisió les dues teories.

La \textbf{teoria de Nivell~I} d'aquest llibre ---els capítols~1 a~18--- és
ETCS: un topos ben apuntat amb objecte dels nombres naturals i elecció, de
la força de la teoria de Zermelo acotada amb elecció. \emph{Tota} la teoria
de l'infinit fins als cardinals hi cap: Cantor, Schröder--Bernstein, bons
ordres, recursió transfinita, Zorn, Hartogs, tipus ordinals, aritmètica
ordinal i cardinal. La raó per la qual hi cap, malgrat que la presentació
clàssica de bona part d'això usa el reemplaçament, és que hem mantingut
totes les construccions \emph{acotades}: cada recursió té codomini fix, cada
suprem de tipus es pren dins un bon ordre ja construït, i la cota de Hartogs
substitueix a tot arreu la recursió \enquote{sobre tots els ordinals}. Aquest
és el contingut de la traducció que el llibre proposa: allà on la teoria
clàssica reuniria una família transfinita amb el reemplaçament, l'estructural
la manté dins un \(\Hart\) i la tracta amb la recursió de codomini fix.

La \textbf{teoria reforçada}, ETCS més la substitució estructural SR
(principi~\ref{sr:esquema}), només fa falta quan la família transfinita
\emph{no es pot acotar}: la successió \(\alpha\mapsto\aleph_\alpha\), els
beths, les jerarquies iterades \(A,P(A),P^2(A),\dots\) reunides en un sol
objecte. SR és la traducció estructural de l'\emph{esquema de reemplaçament}:
allà on ZF diu \enquote{la imatge d'un conjunt per una funció-classe és un
conjunt}, SR diu \enquote{una regla externa uniforme i invariant es
representa per una família interna}. Amb SR afegit, la teoria assoleix la
força de ZFC, i la correspondència amb la teoria clàssica esdevé exacta en el
sentit precís de la bi-interpretabilitat ---no de completesa lògica---, per a
la formulació concreta de l'esquema adoptada (observació~\ref{obs:sr-rigor})
\cite{osius1974,shulman-comparing}.

\begin{observacio}[SR és un esquema, no un axioma; i és opcional]\label{obs:sr-metateoric}
Dues precisions tanquen aquesta comparació. Primera, sobre l'\emph{estatus}:
SR no és un axioma més al costat de l'elecció, sinó un \emph{esquema} de
fórmules, com el reemplaçament clàssic; per això el llibre el presenta com un
principi a part (principi~\ref{sr:esquema}) i el marca separadament a la
caixa de dependències, no barrejat amb els axiomes. Segona, sobre
l'\emph{abast}: SR és estrictament opcional. El volum es pot llegir sencer
fins al capítol~18 sense ell; només el capítol dels alephs
(capítol~\ref{cap:alephs}) el requereix, i ho fa de manera explícita a cada
ús. Un lector que vulgui quedar-se en ETCS pur té, doncs, tota la teoria de
l'infinit numerable, ordinal i cardinal a la seva disposició; qui vulgui els
trams de la successió dels alephs i dels beths com a famílies ha d'adoptar SR, sabent que amb
això puja al nivell de força de ZFC. (El continu \(\mathfrak c=\abs{P(\N)}\) i
la hipòtesi del continu, en canvi, ja són accessibles sense SR, al Nivell~I.)
Aquesta és la línia divisòria exacta entre les dues teories, i el llibre l'ha
fet visible a cada pas en comptes d'amagar-la.
\end{observacio}

\section{Taula comparativa: ZFC, ETCS i aquest llibre}\label{sec:taula-comparativa}

Recollim en tres taules els axiomes de les tres presentacions, i en una
quarta la correspondència entre elles. Per a ZFC seguim la formulació dels
manuals~\cite{jech2003,kunen2011}; per a ETCS, la forma moderna que en donen
McLarty i Shulman~\cite{mclarty-catset,shulman-comparing}, equivalent a
l'original de Lawvere~\cite{lawvere1964,lawvere-etcs-reprint}.

\begin{center}\small
\begin{tabular}{@{}>{\raggedright\arraybackslash}p{0.30\linewidth}>{\raggedright\arraybackslash}p{0.62\linewidth}@{}}
\toprule
\multicolumn{2}{@{}l}{\textbf{ZFC} (primitius: conjunts i \(\in\))}\\
\midrule
Extensionalitat & Dos conjunts amb els mateixos elements són iguals.\\
Fonament & Tot conjunt no buit té un element disjunt d'ell.\\
Separació (esquema) & \(\{x\in A:\varphi(x)\}\) és un conjunt, per a tota fórmula \(\varphi\).\\
Parells & \(\{a,b\}\) és un conjunt.\\
Unió & \(\bigcup A\) és un conjunt.\\
Parts & \(P(A)\) és un conjunt.\\
Infinit & Hi ha un conjunt inductiu.\\
Reemplaçament (esquema) & La imatge d'un conjunt per una funció-classe és un conjunt.\\
Elecció & Tota família de conjunts no buits té una funció d'elecció.\\
\bottomrule
\end{tabular}
\end{center}

\begin{center}\small
\begin{tabular}{@{}>{\raggedright\arraybackslash}p{0.30\linewidth}>{\raggedright\arraybackslash}p{0.62\linewidth}@{}}
\toprule
\multicolumn{2}{@{}l}{\textbf{ETCS} (primitius: objectes, fletxes i composició)}\\
\midrule
Categoria & La composició és associativa i té identitats.\\
Límits finits & Hi ha objecte terminal, productes binaris i equalitzadors.\\
Exponencials & Per a tots \(A,B\) hi ha \(B^A\) amb avaluació universal.\\
Classificador & Hi ha \(\Omega\) que classifica els subobjectes.\\
Nombres naturals & Hi ha un objecte de nombres naturals.\\
Ben apuntat & \(\mathbf1\) separa les fletxes, i \(\mathbf0\not\cong\mathbf1\).\\
Elecció & Tot epimorfisme té secció.\\
\bottomrule
\end{tabular}
\end{center}

\begin{center}\small
\begin{tabular}{@{}>{\raggedright\arraybackslash}p{0.30\linewidth}>{\raggedright\arraybackslash}p{0.62\linewidth}@{}}
\toprule
\multicolumn{2}{@{}l}{\textbf{Aquest llibre} (primitius: aplicacions; \(\Dom\), \(\Cod\), \(\Com\))}\\
\midrule
Estructurals (5) & Domini i codomini, composició, associativitat, neutralitat (capítol~\ref{cap:aplicacions}).\\
Objectes distingits & Terminal, inicial, no-degeneració, separació global (capítol~\ref{cap:aplicacions}).\\
Construccions finites & Sumes, disjunció i estabilitat, productes, representació de les parts per \(\mathbf2\), coequalitzadors (capítols~\ref{cap:parts} i~\ref{cap:valors}).\\
Exactitud & Efectivitat de les relacions d'equivalència (capítol~\ref{cap:relacions}).\\
Exponencials & Existència de \(B^A\) (capítol~\ref{cap:exponencials}).\\
Nombres naturals & Recursió sobre \(\N\) (capítol~\ref{cap:naturals}).\\
Elecció & Tot epimorfisme té secció, i la numerable com a cas particular (capítols~\ref{cap:relacions} i~\ref{cap:infinitud}).\\
SR (esquema, opcional) & Reunió de famílies determinades llevat d'isomorfisme únic (capítol~\ref{cap:alephs}).\\
\bottomrule
\end{tabular}
\end{center}

La taula següent posa les tres presentacions en correspondència. La darrera
columna diu l'estatus de l'axioma del llibre respecte d'ETCS:
\emph{equivalent} (és el mateix axioma, en llenguatge d'aplicacions),
\emph{derivat} (és un teorema d'ETCS que el llibre postula perquè el
necessita abans de tenir les eines per demostrar-lo) o \emph{addicional}
(va més enllà d'ETCS).

\begin{center}\footnotesize
\begin{tabular}{@{}>{\raggedright\arraybackslash\hspace{0pt}}p{0.17\linewidth}>{\raggedright\arraybackslash\hspace{0pt}}p{0.21\linewidth}>{\raggedright\arraybackslash\hspace{0pt}}p{0.16\linewidth}>{\raggedright\arraybackslash\hspace{0pt}}p{0.18\linewidth}>{\raggedright\arraybackslash\hspace{0pt}}p{0.13\linewidth}@{}}
\toprule
Contingut & ZFC & ETCS & Aquest llibre & Estatus\\
\midrule
Composició i identitats & Implícites (funcions com a conjunts de parells) & Categoria & 5 axiomes estructurals & equivalent\\
Terminal & \(\{\varnothing\}\) (Parells) & Límits finits & Axioma & equivalent\\
Inicial & \(\varnothing\) (Separació) & Teorema en un topos & Axioma & derivat\\
Productes & Parells, Parts, Separació & Límits finits & Axioma & equivalent\\
Equalitzadors & Separació & Límits finits & Teorema (cap.~\ref{cap:parts}) & equivalent\\
Sumes; disjunció i estabilitat & Parells, Unió & Teoremes en un topos & Axiomes & derivats\\
Coequalitzadors & Parts, Separació (classes) & Teorema en un topos & Axioma & derivat\\
Efectivitat & Classes d'equivalència & Teorema en un topos & Axioma & derivat\\
Classificador & Separació i Parts (característiques) & Classificador \(\Omega\) & Representació per \(\mathbf2\) & equivalent (\(\Omega\cong\mathbf2\))\\
Exponencials, \(P(A)\) & Parts, Separació & Exponencials & Axioma; \(P(A):=\mathbf2^A\) & equivalent\\
Infinit & Infinit & Nombres naturals & Recursió sobre \(\N\) & equivalent\\
Extensionalitat & Extensionalitat & Ben apuntat & Separació global i no-degeneració & equivalent\\
Elecció & Elecció & Elecció & Elecció & equivalent\\
Separació & Esquema per a tota fórmula & Només fórmules acotades, via \(\Omega\) & Igual que ETCS & equivalent\\
Reemplaçament & Esquema & Absent & SR, opcional & addicional\\
Fonament & Fonament & Sense contrapartida & Sense contrapartida (no hi ha \(\in\) global) & ---\\
\bottomrule
\end{tabular}
\end{center}

\begin{observacio}[Lectura de la taula]\label{obs:lectura-taula}
Tres conclusions. \emph{Primera}: sense SR, els axiomes del llibre són
exactament ETCS. Els que figuren com a \enquote{derivats} ---l'inicial, les
sumes amb la seva disjunció i estabilitat, els coequalitzadors i
l'efectivitat--- són teoremes en qualsevol topos, demostrables a partir
dels límits finits, els exponencials i el classificador (vegeu el volum de
categories~\cite{perello-categories} i~\cite[cap.~IV]{maclane-moerdijk}); el
llibre els postula perquè els necessita als capítols~\ref{cap:parts}
a~\ref{cap:relacions}, abans de disposar del conjunt de parts amb què es
demostrarien. Com a teoria, doncs, el llibre no és més fort que ETCS:
només és redundant. \emph{Segona}: la diferència amb ZFC és doble. Hi ha
dos axiomes de ZFC sense contrapartida ---el \emph{fonament}, que parla
d'una pertinença global que aquí no existeix, i el \emph{reemplaçament},
que aquí és l'esquema opcional SR---, i la \emph{separació} és més feble,
restringida a les condicions que es poden expressar com a parts, és a dir,
amb quantificadors acotats per conjunts ja donats. \emph{Tercera}: que el
classificador sigui el conjunt \(\mathbf2=\mathbf1+\mathbf1\), de dos punts,
és un tret dels conjunts constants: en un topos qualsevol el classificador
\(\Omega\) pot tenir més valors de veritat, i és la combinació de ben
apuntat i elecció la que el redueix a \(\mathbf2\) (Diaconescu: l'elecció
fa booleà el topos). Per això el llibre el pot postular directament.
\end{observacio}

\begin{observacio}[Què és demostrat i què és citat]\label{obs:taula-citat}
La taula barreja dues menes d'afirmacions, que cal no confondre. Les
columnes \enquote{ZFC} i \enquote{Aquest llibre} i la correspondència entre
elles en el pla del \emph{contingut} (quin principi fa la feina de quin) són
comparacions informals que el lector pot comprovar amb el llibre a la mà.
En canvi, les afirmacions sobre \emph{teories} ---que els axiomes
\enquote{derivats} són teoremes de tot topos, que ETCS és equivalent a BZC, i
que ETCS\({}+{}\)SR és bi-interpretable amb ZFC--- són \emph{metateoremes
citats} de la literatura~\cite{osius1974,mclarty-catset,shulman-comparing,maclane-moerdijk},
no resultats demostrats en aquest volum, i la darrera depèn de la
formulació concreta de SR (observació~\ref{obs:sr-rigor}). El que el llibre
sí que demostra és la direcció directa: que tots els seus teoremes són
conseqüència dels axiomes de la taula.
\end{observacio}

\begin{observacio}[Força de les tres teories]\label{obs:forca-taula}
En força, l'escala és: \(\text{Nivell~I}=\text{ETCS}\equiv\mathrm{BZC}\)
(Zermelo amb elecció i separació acotada) \(<\) \(\mathrm{ZFC}\equiv
\text{ETCS}+\text{SR}\), on \(\equiv\) indica bi-interpretabilitat i \(<\) que
la teoria de la dreta demostra la consistència de la de l'esquerra
(observació~\ref{obs:forca-relativa};
\cite{osius1974,mclarty-catset,shulman-comparing}). La mancança de
reemplaçament no es nota en la matemàtica ordinària, i la part transfinita
del llibre s'ha mantingut al Nivell~I gràcies a les cotes de Hartogs; només
la jerarquia sencera dels alephs i dels beths l'exigeix.
\end{observacio}

\section{Límits del desenvolupament}\label{sec:limits}

Tanquem amb els límits, que un capítol de balanç honest no pot ometre.
Són de tres menes: de mida, de consistència i de força expressiva.

\subsection{Mida}\label{subsec:mida}

Els axiomes no diuen res sobre \emph{quants} conjunts hi ha ni sobre com
n'és de gran l'univers. Garanteixen que hi ha un conjunt infinit
(\(\N\)), i que a partir de qualsevol conjunt se'n poden formar el de parts
i els exponencials ---de manera que hi ha conjunts arbitràriament grans en
el sentit de Cantor---, però no postulen cap cardinal inaccessible ni cap
tancament sota famílies de mida arbitrària. En particular, sense afegir SR,
no hi ha reemplaçament: no es pot, en general, formar la unió d'una família
de conjunts indexada per un conjunt quan la família ve donada per una
\emph{fórmula} i no per una aplicació ja existent ---és exactament la
distinció entre regla externa i família interna que vam fer al
capítol~\ref{cap:families}. Aquesta és la limitació de mida més important de
la teoria de Nivell~I, i és exactament la que SR (principi~\ref{sr:esquema})
aixeca, permetent reunir les famílies transfinites no acotades del
capítol~\ref{cap:alephs}. Amb SR, la teoria té la força del reemplaçament;
sense ell, es queda ---còmodament--- en tota la matemàtica que no reuneix
col·leccions transfinites no acotades.

\subsection{Consistència relativa}\label{subsec:consistencia}

Com qualsevol teoria que interpreti l'aritmètica, la nostra no pot
demostrar la seva pròpia consistència, pel segon teorema de Gödel. El que
sí que tenim són resultats de consistència \emph{relativa}: la teoria és
consistent si ho és la teoria de Zermelo amb elecció, i recíprocament, per
la interpretació mútua esmentada. Dit d'una altra manera, no demanem al
lector cap acte de fe que no demani ja la teoria clàssica: qui accepti la
consistència de \(\mathrm{ZFC}\) ---o fins i tot només la de \(Z\) amb
elecció--- accepta la d'aquest llibre. No n'obtenim, doncs, cap garantia
absoluta de consistència; només la seguretat que no hem introduït cap risc
nou respecte del que la matemàtica ja assumeix.

\subsection{Força expressiva}\label{subsec:expressiva}

Finalment, hi ha coses que la teoria, tal com és, expressa amb dificultat
o no expressa. Les propietats que \emph{no} són invariants per isomorfisme
---els fets de codificació interna, la pertinença creuada entre conjunts
sense relació--- no s'hi poden ni tan sols formular, i això, com hem
argumentat, és volgut. Pel que fa als arguments transfinits, la situació és
més matisada del que la primera versió d'aquest capítol suggeria: la teoria
de l'ordre, els tipus ordinals i la recursió transfinita \emph{sí} que hi
són, sencers, mentre es mantinguin \emph{acotats} ---la recursió amb codomini
fix (capítol~\ref{cap:recursio}), els suprems dins un bon ordre, la cota de
Hartogs. El que queda fora de la teoria de Nivell~I, sense SR, és
exactament la reunió de col·leccions transfinites \emph{no acotades}: els
trams de la successió dels \(\aleph_\alpha\) i \(\beth_\alpha\), la jerarquia
acumulativa \(V_\alpha\) per a ordinals grans, o la recursió al llarg de
qualsevol ordinal fix per molt llarg que sigui. Tot això és accessible amb SR
(capítol~\ref{cap:alephs}), sempre tram a tram: ni amb SR no hi ha cap
objecte indexat per tots els ordinals;
i sense SR es queda fora ---no per una mancança del llenguatge, sinó per la
mateixa raó que ho queda a ETCS: manca el reemplaçament. Res d'això no és
una objecció: és la delimitació natural d'una teoria pensada per fonamentar
la matemàtica ordinària amb el mínim de compromisos, i que fa explícit,
mitjançant SR, el punt on cal pagar-ne un de més fort.

\begin{observacio}[Balanç del capítol i del llibre]\label{obs:balanc-metateoric}
Aquest capítol ha respost les tres preguntes amb què s'obria. \emph{És
realment una teoria de conjunts?} Sí: és una axiomatització d'ETCS,
autònoma, en un llenguatge d'aplicacions
(observació~\ref{obs:tesi-etcs}), i la categoria dels conjunts clàssics
n'és un model (secció~\ref{sec:conjunts-petits}). \emph{Què pot expressar?}
Tota la matemàtica que viu llevat d'isomorfisme ---que és, a la pràctica,
tota la matemàtica ordinària---, amb la força de la teoria de Zermelo amb
elecció; i, dins d'aquesta, tota la teoria clàssica de l'infinit ---Cantor,
bons ordres, recursió transfinita, ordinals i cardinals--- \emph{acotada},
sense cap principi nou (secció~\ref{sec:transfinit-balanc}). Només la reunió
de col·leccions transfinites no acotades ---els trams de la successió dels alephs i
dels beths--- demana la substitució estructural SR, que porta la teoria a la
força de ZFC (secció~\ref{sec:etcs-sr}); el continu mateix
\(\mathfrak c=\abs{P(\N)}\) i la hipòtesi del continu, en canvi, són ja de
Nivell~I. \emph{Quina relació té amb la teoria clàssica?}
És equiconsistent i mútuament interpretable amb \(Z\) acotada amb elecció,
una mica més feble que ZFC, i coincideix amb ZFC quan s'hi afegeix SR ---la
traducció estructural del reemplaçament (observació~\ref{obs:forca-relativa},
secció~\ref{sec:etcs-sr}); la diferència essencial és que substitueix la
pertinença primitiva i global per una de derivada i local
(secció~\ref{sec:pertinenca}), i que tracta ordinals i cardinals com a
\emph{tipus} llevat d'isomorfisme en comptes de com a conjunts privilegiats
(secció~\ref{sec:vonneumann}).

Amb això es tanca el recorregut començat a la introducció. Vam prometre
construir els conjunts sense posar els elements al centre, prenent les
aplicacions com a únics primitius i recuperant-ho tot ---parts,
pertinença, productes, quocients, funcions, nombres i cardinalitat--- a
partir del domini, el codomini i la composició. El recorregut ha arribat,
partint d'aquí, fins a la teoria clàssica de l'infinit sencera: mides i
formes d'ordre, cada una amb la seva aritmètica, els dos grans teoremes
diagonals ---Cantor i Burali--Forti--- que impedeixen una mida o una
longitud màximes, i, amb SR, la jerarquia dels alephs i la formulació precisa
de la hipòtesi del continu. Tot dit en un llenguatge que posa en primer pla
les relacions entre conjunts en comptes de la seva composició interna. Qui
vulgui anar més enllà ---grans cardinals, forcing, o la traducció detallada
cap a l'univers de von Neumann i tornada--- té, en la
bibliografia~\cite{mclarty-catset,shulman-comparing,lawvere-rosebrugh}, els fils
per continuar.
\end{observacio}
